%% CAPITULO 4 -- ANALISE E DISCUSSAO

\section{Análise e discussão dos dados}

Esta versão do trabalho é um esqueleto: ainda não há dados próprios, e a análise se restringe ao que as obras de referência relatam em seus resumos. Os números citados a seguir pertencem aos estudos originais e não podem ser comparados diretamente entre si, pois foram obtidos em contextos e com métricas distintos.

\citeonline{majumder2025} relatam que o fluxo com modelo de linguagem filtrou 80,5\% das denúncias irrelevantes e elevou a acurácia de roteamento de 31,8\% para 82,4\% em relação ao sistema vigente. \citeonline{buryakov2024} obtiveram F1 médio de 0,61 sobre 12 categorias, e \citeonline{prencipe2026} relatam redução de 70\% no tempo de preparo de respostas.

Em comum, os três estudos tratam de demandas dirigidas a órgãos públicos, mas apenas o último descreve explicitamente a supervisão humana como parte do fluxo. Essa diferença motiva a discussão posterior sobre como combinar classificação automática, rastreabilidade e revisão por servidores.
